\documentclass[12pt, english]{article}
%%%%%%%%%%%%%%%%%%%%%%%%%%%%%%%%%%%%%%%%%%%%%%%%%%%%%%%%

% margin %%%%%%%%%%%%%%%%%%%%%%
\usepackage{geometry}
 \geometry{
 a4paper,
 left=1in,
 top=0.9in,
 right=1in,
 bottom=2in,
 }
%%%%%%%%%%%%%%%%%%%%%%%%%%%%%%%%%%%%%%%%%%%%%%%%%%%%%%%%

% page no. %%%%%%%%%%%%%%%%%%%%%%
\usepackage{fancyhdr}
\pagestyle{fancy}
\fancyhf{}
\fancyfoot[C]{\textit{\thepage}}
\setlength{\headheight}{65pt}
%%%%%%%%%%%%%%%%%%%%%%%%%%%%%%%%%%%%%%%%%%%%%%%%%%%%%%%%

% idk-needed %%%%%%%%%%%%%%%%%%%%%%
\usepackage{amsmath,amsthm,amssymb,amsfonts, enumitem, color, comment, graphicx, environ, lipsum, accsupp}
%%%%%%%%%%%%%%%%%%%%%%%%%%%%%%%%%%%%%%%%%%%%%%%%%%%%%%%%

% code and color %%%%%%%%%%%%%%%%%%%%%%
\usepackage{spverbatim}
\usepackage{xcolor}
\usepackage{listings}
\definecolor{bluegrey}{RGB}{0, 0, 200}
\lstset{
        framexleftmargin=0pt,
        frame=shadowbox,
        rulesepcolor=\color{gray},
        xleftmargin=30pt,
        xrightmargin=30pt,
        breaklines=true,
        basicstyle=\ttfamily,
        showstringspaces=false,
        escapeinside={(*@}{@*)},
        }
\def\speciallstcolor{\begingroup\color{blue}}
\def\endspeciallstcolor{\endgroup}
%%%%%%%%%%%%%%%%%%%%%%%%%%%%%%%%%%%%%%%%%%%%%%%%%%%%%%%%

% heading %%%%%%%%%%%%%%%%%%%%%%
\renewcommand\thesection{\color{blue}\arabic{section}}
\newcommand{\mysection}[2]{\setcounter{section}{#1}\addtocounter{section}{-1}\section{#2}}
%%%%%%%%%%%%%%%%%%%%%%%%%%%%%%%%%%%%%%%%%%%%%%%%%%%%%%%%

% both side indent %%%%%%%%%%%%%%%%%%%%%%
\newenvironment{myindent}
{\par\leftskip25pt\relax\rightskip25pt\relax}
{\par\leftskip0cm\relax\rightskip0cm\relax}
%%%%%%%%%%%%%%%%%%%%%%%%%%%%%%%%%%%%%%%%%%%%%%%%%%%%%%%%

% type ^ and \ %%%%%%%%%%%%%%%%%%%%%%
\newcommand{\power}{$^\wedge$}
\newcommand{\backs}{\textbackslash} 
%%%%%%%%%%%%%%%%%%%%%%%%%%%%%%%%%%%%%%%%%%%%%%%%%%%%%%%%

% document info %%%%%%%%%%%%%%%%%%%%%%
\lhead{Arav Sarma \\ 15193845}  %replace with your name
\rhead{Homework \textcolor{blue}{6} \\ 2025 Speech Processing}
%%%%%%%%%%%%%%%%%%%%%%%%%%%%%%%%%%%%%%%%%%%%%%%%%%%%%%%%

% begin %%%%%%%%%%%%%%%%%%%%%%
\begin{document}
%wwwwwwwwwwwwwwwwwwwwwwwwwwwwwwwwwwwwwwwwwwwwwwwwwwwwwwwww

% section _______________________________________________
\mysection{6}{Pitch}
\vspace{10pt}
% -__-__-__-__-__-__-__-__-__-__-__-__-__-__-__-__-__-__-

% //////////////////////////////////////////////////////
\subsection*{\emph{2.} Automated pitch measurement} \vspace{15pt}

\begin{itemize}
    \item \textbf{Script} \vspace{5pt}
\begin{lstlisting}[xleftmargin=0pt,]
textgridFiles$# = fileNames$# ("sounds/*.TextGrid")
writeInfoLine: "Durations and pitches:", newline$
for fileNumber to size (textgridFiles$#)
	textgridFile$ = textgridFiles$# [fileNumber]
	textgrid = Read from file: "sounds/" + textgridFile$
	numberOfIntervals = Get number of intervals: 1
	soundFile$ = textgridFile$ - ".TextGrid" + ".wav"
	sound = Read from file: "sounds/" + soundFile$
	pitchObject = To Pitch: (*@\aftergroup\speciallstcolor@*)0, 75, 600(*@\aftergroup\endspeciallstcolor@*)
	
	for interval to numberOfIntervals
		selectObject: (*@\aftergroup\speciallstcolor@*)textgrid(*@\aftergroup\endspeciallstcolor@*)
		text$ = Get label of interval: 1, interval
		if text$ = "v"
			start = Get start time of interval: 1, interval
			end = Get end time of interval: 1, interval
			duration = end - start
			selectObject: pitchObject
			pitch = Get mean: (*@\aftergroup\speciallstcolor@*)0, 0, "Hertz"(*@\aftergroup\endspeciallstcolor@*)
			appendInfoLine: "File ", textgridFile$ - ".TextGrid", ": ", "duration = ", fixed$ (duration, 3), " seconds,", " pitch = ", fixed$ (pitch, 3), " Hz."
		endif
	endfor
	removeObject: (*@\aftergroup\speciallstcolor@*)textgrid, sound, pitchObject(*@\aftergroup\endspeciallstcolor@*)
endfor
\end{lstlisting} \vspace{10pt}

    \item \textbf{Output} \vspace{5pt}
\begin{lstlisting}[xleftmargin=0pt]
Durations and pitches:

File ki-ba1: duration = 0.060 seconds, pitch = 147.511 Hz.
File ki-ba2: duration = 0.075 seconds, pitch = 141.793 Hz.
File ki-ba3: duration = 0.090 seconds, pitch = 146.283 Hz.
File kiba1: duration = 0.032 seconds, pitch = 131.785 Hz.
File kiba2: duration = 0.039 seconds, pitch = 131.806 Hz.
File kiba3: duration = 0.055 seconds, pitch = 131.514 Hz.
\end{lstlisting} \vspace{10pt}

\end{itemize}

% //////////////////////////////////////////////////////

% \\\\\\\\\\\\\\\\\\\\\\\\\\\\\\\\\\\\\\\\\\\\\\\\\\\\\\\\\\
\subsection*{\emph{3.} Pitch manipulation}

\begin{itemize}
    \item \textbf{Note} \par
   The synthesised question sound did not sound exactly like a question to me even after I matched the countour to a question recording of the same sentence I made. I suspect this is because of duration differences between declarative and interrogative prosody.
\end{itemize}

% \\\\\\\\\\\\\\\\\\\\\\\\\\\\\\\\\\\\\\\\\\\\\\\\\\\\\\\\\\

% ><><><><><><><><><><><><><><><><><><><><><><><><><><><><><
\subsection*{\emph{4.} Pitch in the acoustic continua of your experiment}
\vspace{10pt}

\begin{itemize}
    \item \textbf{General research question:} \par
    \textcolor{gray}{Location of phoneme boundary as a function of orthographic context (and exopure/familiarity with orthography of certain language(s))} \par
    \textbf{Location of phoneme boundary as a function of lexical context (word-nonword)}
    \item \textbf{\textcolor{gray}{Addition:}} \par
    \textcolor{gray}{L2 lexical knowledge context}
    \item \textbf{Specific research question:} \par
    \textcolor{gray}{My idea is to take both orthography and non-native lexical knowledge as a coupled effect on phonological categorisation. Preliminary thoughts on operationalising this question: looking at phonemic boundary between voiced-plain-aspirated, aspirated-fricative continua for L2 English speakers from Assam (or L1 Assamese speaking). Onset stop contrast on L2 variety is through voicing unlike aspiration in L1 English for Assamese (voicing+aspiration for Boro), Orthograhpic <b>, <d>, <g> are percieved as voiced like other Indian Englishes. This perception gets carried over even when exposed to L1 English. Questions: How much closure devoicing / afterburst duration can be blocked off during phoneme categorisation for a voicing / plain percept (how strong would lexical effect be on a word like `ban' / `pan', and also orthographic; eg: when shown <pan> / <phan or fan>)?} \par
    How does presence/absence of a lexical item effect categorisation of aspirated-breathy laryngeal contrast? Unlike the voicing or aspirated contrasts where boundary can be tested with VOT continua, the aspirated-breathy contrast would require manipulation of cues like turbulence quality and also pitch on following vowel (as the breathy and aspirated categories has been observed to effect pitch across languages; tonogenesis in Rohingya, Punjabi, tone-depressing effect in Southeastern Bantu, adaptation of indic loans with breathy in Thai). The study will look at lexical effect across this contrast for L1 Assamese speakers from Guwahati.       
    \item \textbf{Participants} \par
    Recruited: \char`\~ 7 \\
    Completed: 0 \par
    Participants are going to be tested online using Zoom online meetings, and recruited through author's contacts, and likely by reaching out to Assam/Northeast related online forums like on reddit. The plan is to have non-native English speakers ideally all from Guwahati with same L1, this is to neutralise individual factors like regional variation, exposure to English. A different experiment setup can be by using a qualtrics form along with online meet, so the participants have more autonomy with the test and excess to stimuli on their device while also making it easier to see if they are in a quiet environment. \par
    \item \textbf{Continua}
    Acoustic continua from [bha] to [pha]: 1st context is with following syllable [ŋa] (word to nonword). 2nd context is with following syllable [ki] (nonword to word). \\
    Continua from [dh] to [th]ː 1st context is with nucleus [ue] (word to nonword). 2nd context is with nucleus [io] (nonword to word). \par
    \item \textbf{Duration} \par
    The breathy-aspirated contrast can involve duration of voicing, however the breathy phonation is not really the same as voicing of a vowel (vocal folds in breathy are not as / or not completely close together as in modal voice). Therefore it is amount of voicing (closeness of vocal folds) during the afterburst which would be contrastive cue in terms of voicing. The other place where duration might be important is onset of vowel (duration of afterbust). However this will not be a cue used for this study's continua as it has not been investigated if it is involved in breathy-aspirated contrast and because onset of vowel is not easy to determine, especially if some amount of voicing is present before the vowe which can be the case for breathy phonemes.  \par
    \item \textbf{Pitch} \par
    Pitch is a cue which will be used to form continua. The vowel portion of the sound will be manipulated for gradient (mean) pitch. To synthesise different pitches, first the target pitch points in the manipulation window are selected, then they are scaled up/down to desired mean pitch via \textbf{Multiply pitch frequencies...}. \par
    Points to consider:
    \begin{itemize}
        \item How is pitch gradient incorporated with degree of voicing gradient: congruent + non-congruent?
        \item Should a variety of vocal tract size effect on perceived pitch be present on the stimuli? 
    \end{itemize}
\end{itemize}
% ><><><><><><><><><><><><><><><><><><><><><><><><><><><><><

%wwwwwwwwwwwwwwwwwwwwwwwwwwwwwwwwwwwwwwwwwwwwwwwwwwwwwwwww
\end{document}